% Auto-generated by data/nordic_summaries/combine_nordic_summaries.py
% Source: data/nordic_summaries/nordic_summaries_combined.csv
% Do not edit by hand -- rerun the script instead.
%
% Requires: longtable, booktabs, amssymb (\checkmark), natbib (\citet).
%
% Basis for each key-pole flag:
%   Long Range Dykes: Buchan (2013) Table 1. NOTE: flagged to stay faithful to his call, but this is the weakest key pole in the compilation and warrants discussion in the text -- his pole rests on 5 dykes with A95 18 deg, and the pole recreated here has A95 20.8 deg, so secular variation is unlikely to be averaged. It also sets the effective floor on how permissive his quality criterion is; baked contact test
%   Franklin event grand mean: Buchan (2013) Table 1; regional consistency test
%   Gunbarrel LIP: Buchan (2013) Table 1; regional consistency test
%   Jacobsville Formation: two positive intraformational conglomerate tests within the formation (Agate Falls, Dover Creek); inclination shallowing displaces the pole by 3.5 deg. NOTE: the age is equivocal against the +/-15 Myr standard, resting on a detrital-zircon maximum depositional age of 992.51 +/- 0.64 Ma rather than a date on the unit, and the reversal test fails; intraformational conglomerate test
%   Lower Freda Formation: the intraformational conglomerate test of Swanson-Hysell, Fairchild & Slotznick (2019b) is on intraclasts within this formation; inclination shallowing displaces the pole by 0.1 deg; intraformational conglomerate test
%   Nonesuch Shale: intraformational conglomerate test on fluvial intraclasts in the conformably overlying lower Freda Formation (Swanson-Hysell, Fairchild & Slotznick, 2019b), of the same lithofacies as, and in close stratigraphic proximity to, the fluvial rocks that carry this pole; inclination shallowing is quantified and displaces the pole by only 3.7 deg; intraformational conglomerate test
%   Cardenas Basalts: regional consistency with the coeval Lake Shore Traps and Michipicoten Island poles ~2400 km away; regional consistency test
%   Michipicoten Island Formation: regional consistency with the coeval Cardenas Basalts pole; regional consistency test
%   Lake Shore Traps: Buchan (2013) Table 1; baked contact test
%   Schroeder Lutsen Basalts: regional consistency with the coeval Lake Shore Traps and Portage Lake Volcanics poles; regional consistency test
%   Portage Lake Volcanics: Buchan (2013) Table 1; secular variation correlation test
%   Uppermost Mamainse Point volcanics -N: remanence direction correlation between the Mamainse Point polarity zones; regional consistency with the coeval Portage Lake Volcanics pole; remanence direction correlation test, regional consistency test
%   North Shore Volcanic Group -N (combined): regional consistency with the coeval Portage Lake Volcanics and uppermost Mamainse Point poles; regional consistency test
%   Mamainse Point volcanics -C (lower N, upper R): Great Conglomerate above the zone carries clasts derived from it (Swanson-Hysell et al., 2009); intraformational conglomerate test
%   Lower Mamainse Point volcanics -R2: conglomerate above the zone carries clasts derived from it (Swanson-Hysell et al., 2009); intraformational conglomerate test
%   Osler Volcanic Group reverse upper: Buchan (2013) Table 1; fold test contemporaneous with emplacement
%   Coldwell Complex: regional consistency with the coeval upper Osler and Nipigon (Logan) sill poles; regional consistency test
%   Osler Volcanic Group reverse middle: regional consistency with the coeval Nipigon (Logan) sill and lowermost Mamainse Point poles; regional consistency test
%   Nipigon sills and lavas: Buchan (2013) Table 1 (Logan sills); baked contact test
%   Osler Volcanic Group reverse lower: regional consistency with the coeval Nipigon (Logan) sill pole; regional consistency test
%   Lowermost Mamainse Point volcanics -R1: remanence direction correlation between the Mamainse Point polarity zones; regional consistency with the coeval Nipigon (Logan) sill pole; remanence direction correlation test, regional consistency test
%   Abitibi Dykes: Buchan (2013) Table 1; baked contact test
%   NW Ontario Lamprophyre Dykes: positive baked-contact test on dykes D1 and MM2 (Piispa et al., 2018); baked contact test
%   Sudbury Dike Swarm: Buchan (2013) Table 1; baked contact profile test
%   Mackenzie dykes recompiled: Buchan (2013) Table 1; baked contact test
%   Midsommersoe Dolerites: Buchan (2013) Table 1; baked contact test
%   Victoria Fjord dolerite dykes: component of the Buchan (2013) Table 1 Midsommerso grand mean; carries the positive baked-contact test he cites for it. NOTE: these dykes have no direct radiometric date -- their age rests on the antipodality of their directions to the dated Midsommerso dolerites, which Buchan accepted in folding them into one 1382 +/- 2 Ma pole; baked contact test
%   Zig-Zag Dal Basalts: component of the Buchan (2013) Table 1 Midsommerso grand mean; same 1382 +/- 2 Ma event as the baked-contact-tested Victoria Fjord dykes; baked contact test, regional consistency test
%   Mean Rocky Mountain intrusions: Buchan (2013) Table 1 (Laramie complex and Sherman Granite) -- regional consistency; regional consistency test
%   Mistastin Batholith: Buchan (2013) Table 1 -- regional consistency with the ca. 1.46-1.42 Ga suite, not the impact-crater test; regional consistency test
%   St. Francois Mountains Acidic Rocks: Buchan (2013) Table 1; fold test contemporaneous with emplacement, regional consistency test
%   Michikamau Intrusion Combined: Buchan (2013) Table 1; baked contact test
%   Western Channel diabase: Buchan (2013) Table 1; secular variation correlation test, remanence direction correlation test
%   Cleaver Dykes: Buchan (2013) Table 1; baked contact test
%   Wharton Group: robustly positive, definitively intraformational conglomerate test on rhyolite clasts within the Pitz Formation section at McRae Lake (Raub et al., 2026); the pole is dated directly by a 1756.4 +/- 0.9 Ma U-Pb zircon age on one of its contributing flows. Raub et al. reach the same conclusion independently, writing that the pole "constitutes a key pole by the criteria outlined in Buchan et al. (2000) and Buchan (2013)"; intraformational conglomerate test
%   NE trending ECMB Diabase Dykes: regional consistency with the Wharton Group pole 2027 km away and the Cleaver dykes pole 2659 km away, both of which are key on a direct field test; three discrete igneous units spanning 1779-1740 Ma across the craton interior. NOTE: this application of the regional consistency test is provisional and is the least settled flag in the compilation. Apparent polar wander is slow across this interval -- Raub et al. (2026) call it "a short arc" -- so agreement between poles tens of Myr apart is only weakly diagnostic of primary remanence, unlike in the rapidly moving Midcontinent Rift where coeval poles a few Myr apart routinely fail to overlap. The agreement with Wharton is marginal (15.6 deg against a 15.9 deg summed A95); the agreement with Cleaver is comfortable (8.5 deg against 11.5 deg). The pole has no primary field test of its own, only an inverse baked contact against a ca. 1096 Ma dyke. Worth putting to Buchan directly before publication; regional consistency test
%
% CAVEATS -- flags carried for fidelity to Buchan (2013) that warrant
%           discussion in the text:
%   Long Range Dykes: flagged to stay faithful to his call, but this is the weakest key pole in the compilation and warrants discussion in the text -- his pole rests on 5 dykes with A95 18 deg, and the pole recreated here has A95 20.8 deg, so secular variation is unlikely to be averaged. It also sets the effective floor on how permissive his quality criterion is; baked contact test
%   Jacobsville Formation: the age is equivocal against the +/-15 Myr standard, resting on a detrital-zircon maximum depositional age of 992.51 +/- 0.64 Ma rather than a date on the unit, and the reversal test fails; intraformational conglomerate test
%   Victoria Fjord dolerite dykes: these dykes have no direct radiometric date -- their age rests on the antipodality of their directions to the dated Midsommerso dolerites, which Buchan accepted in folding them into one 1382 +/- 2 Ma pole; baked contact test
%   NE trending ECMB Diabase Dykes: this application of the regional consistency test is provisional and is the least settled flag in the compilation. Apparent polar wander is slow across this interval -- Raub et al. (2026) call it "a short arc" -- so agreement between poles tens of Myr apart is only weakly diagnostic of primary remanence, unlike in the rapidly moving Midcontinent Rift where coeval poles a few Myr apart routinely fail to overlap. The agreement with Wharton is marginal (15.6 deg against a 15.9 deg summed A95); the agreement with Cleaver is comfortable (8.5 deg against 11.5 deg). The pole has no primary field test of its own, only an inverse baked contact against a ca. 1096 Ma dyke. Worth putting to Buchan directly before publication; regional consistency test
\begingroup
\footnotesize
\setlength{\tabcolsep}{3pt}
\begin{longtable}{p{1.8cm}p{3.1cm}p{1.5cm}cccrrrrrp{1.45cm}p{2.6cm}}
\caption{Compiled paleomagnetic poles for Laurentia spanning ca.\ 1779--565~Ma, beginning after the amalgamation of Laurentia and running through the remainder of the Proterozoic, listed from youngest to oldest. Each pole is recalculated from site-level data in the accompanying compilation. ``Age (Ma)'' gives the nominal age of magnetization with its lower and upper bounds beneath. ``Rating'' is the Nordic reliability grade (A/B). ``Key pole'' ($\checkmark$) marks poles meeting the key-pole criteria of \citet{Buchan2013a}: a precisely determined age (U--Pb, or occasionally Ar--Ar, typically within $\pm$10~Myr), a good-quality determination in which the primary remanence is isolated by stepwise demagnetization and secular variation is largely averaged, and a positive field test establishing that the remanence is primary rather than solely older than a subsequent geologic event. This flag is assessed independently of the Nordic grade and of the R-criteria scores of \citet{Meert2020a}, which are reported in the accompanying compilation; a blank marks a pole assessed against these criteria and found not to meet them, and a dash marks one not yet assessed. Site and pole longitudes are in degrees east ($0$--$360^{\circ}$), and ``Duluth paleolat'' is the paleolatitude of Duluth, Minnesota ($46.8^{\circ}$N, $267.9^{\circ}$E) implied by the pole. Site and pole positions are present-day; the Greenland and Scotland poles are rotated back into Laurentia coordinates before their Duluth paleolatitude is computed, so that column is comparable across the whole table.}\label{tab:poles} \\
\toprule
\textbf{Terrane} & \textbf{Unit} & \textbf{Age (Ma)} & \textbf{Rating} & \textbf{R1--R7} & \textbf{Key pole} & \textbf{Site lon} & \textbf{Site lat} & \textbf{Plon} & \textbf{Plat} & \textbf{A95} & \textbf{Duluth}\newline\textbf{paleolat} & \textbf{Pole reference} \tabularnewline
\midrule
\endfirsthead
\toprule
\textbf{Terrane} & \textbf{Unit} & \textbf{Age (Ma)} & \textbf{Rating} & \textbf{R1--R7} & \textbf{Key pole} & \textbf{Site lon} & \textbf{Site lat} & \textbf{Plon} & \textbf{Plat} & \textbf{A95} & \textbf{Duluth}\newline\textbf{paleolat} & \textbf{Pole reference} \tabularnewline
\midrule
\endhead
\midrule
\multicolumn{13}{r}{\textit{continued on next page}} \tabularnewline
\endfoot
\bottomrule
\endlastfoot
\raggedright Laurentia & \raggedright Sept-Iles Layered Intrusion & \raggedright 565\newline{\scriptsize 561--569} & B & 1111000 &  & 293.5 & 50.2 & 321.2 & $-$19.7 & 6.5 & \hfill 8.0 & \raggedright \citet{Tanczyk1987a} \tabularnewline
\raggedright Laurentia & \raggedright Catoctin Basalts & \raggedright 568\newline{\scriptsize 555--577} & B & 1011111 &  & 281.8 & 38.5 & 306.9 & 42.4 & 16.5 & \hfill 62.2 & \raggedright \citet{Meert1994a} \tabularnewline
\raggedright Laurentia & \raggedright Callander Alkaline Complex & \raggedright 575\newline{\scriptsize 570--580} & B & 0111101 &  & 280.6 & 46.2 & 300.8 & 46.5 & 5.8 & \hfill 67.6 & \raggedright \citet{Symons1991a} \tabularnewline
\raggedright Laurentia & \raggedright Baie des Moutons complex A & \raggedright 583\newline{\scriptsize 581--585} & B & 1110101 &  & 301.0 & 50.8 & 331.9 & 42.9 & 11.5 & \hfill 45.7 & \raggedright \citet{McCausland2011a} \tabularnewline
\raggedright Laurentia & \raggedright Baie des Moutons complex B & \raggedright 583\newline{\scriptsize 581--585} & B & 1010100 &  & 301.0 & 50.8 & 321.9 & $-$33.9 & 15.7 & \hfill $-$4.2 & \raggedright \citet{McCausland2011a} \tabularnewline
\raggedright Laurentia & \raggedright Long Range Dykes & \raggedright 615\newline{\scriptsize 613--617} & B & 1011101 & $\checkmark$ & 303.3 & 53.7 & 165.4 & $-$12.4 & 20.8 & \hfill $-$17.5 & \raggedright \citet{Murthy1992a} \tabularnewline
\raggedright Laurentia & \raggedright Franklin event grand mean & \raggedright 719\newline{\scriptsize 718--720} & A & 1111111 & $\checkmark$ & 275.4 & 73.0 & 162.1 & 6.7 & 3.0 & \hfill $-$5.7 & \raggedright \citet{Denyszyn2009b} \tabularnewline
\raggedright Laurentia & \raggedright Chuar Group (combined) & \raggedright 755\newline{\scriptsize 748--761} & A & 1111110 &  & 248.1 & 36.2 & 163.6 & 7.5 & 4.2 & \hfill $-$4.2 & \raggedright \citet{Weil2004a,Eyster2020a} \tabularnewline
\raggedright Laurentia & \raggedright Uinta Mountain Group & \raggedright 759\newline{\scriptsize 743--767} & B & 1110100 &  & 250.7 & 40.8 & 160.6 & 1.9 & 2.1 & \hfill $-$10.3 & \raggedright \citet{Weil2006b} \tabularnewline
\raggedright Laurentia & \raggedright Gunbarrel LIP & \raggedright 780\newline{\scriptsize 776--782} & A & 1111101 & $\checkmark$ & 250.1 & 45.0 & 151.5 & 3.2 & 8.0 & \hfill $-$15.3 & \raggedright \citet{Ding2025a} \tabularnewline
\raggedright Laurentia & \raggedright Adirondack metamorphic anorthosite & \raggedright 887\newline{\scriptsize 864--910} & B & 0111101 &  & 286.1 & 44.2 & 143.4 & $-$25.2 & 12.9 & \hfill $-$41.4 & \raggedright \citet{Brown2012a} \tabularnewline
\raggedright Laurentia-Scotland & \raggedright Torridon Group & \raggedright 975\newline{\scriptsize 950--1050} & B & 0111001 &  & 354.3 & 57.9 & 220.9 & $-$17.7 & 7.1 & \hfill $-$8.6 & \raggedright \citet{Evans2021a} \tabularnewline
\raggedright Laurentia & \raggedright Jacobsville Formation & \raggedright 990\newline{\scriptsize 980--1000} & A & 1111101 & $\checkmark$ & 271.8 & 47.3 & 186.2 & $-$14.2 & 2.7 & \hfill $-$4.8 & \raggedright \citet{Zhang2024a} \tabularnewline
\raggedright Laurentia & \raggedright Upper Freda Formation & \raggedright 1045\newline{\scriptsize 1032--1063} & B & 0111101 &  & 271.0 & 47.0 & 169.8 & 4.7 & 2.1 & \hfill $-$2.1 & \raggedright \citet{Fuentes2025a} \tabularnewline
\raggedright Laurentia & \raggedright Lower Freda Formation & \raggedright 1075\newline{\scriptsize 1063--1085} & B & 0111001 & $\checkmark$ & 270.4 & 46.8 & 179.7 & 1.2 & 2.4 & \hfill 2.1 & \raggedright \citet{Henry1977a} \tabularnewline
\raggedright Laurentia & \raggedright Nonesuch Shale & \raggedright 1077\newline{\scriptsize 1066--1086} & A & 1111101 & $\checkmark$ & 269.5 & 46.5 & 180.4 & 3.9 & 1.7 & \hfill 4.6 & \raggedright \citet{Slotznick2024a} \tabularnewline
\raggedright Laurentia & \raggedright Cardenas Basalts & \raggedright 1082\newline{\scriptsize 1081--1083} & A & 1110101 & $\checkmark$ & 248.1 & 36.3 & 183.9 & 15.9 & 7.4 & \hfill 15.6 & \raggedright \citet{Zhang2024b} \tabularnewline
\raggedright Laurentia & \raggedright Michipicoten Island Formation & \raggedright 1084\newline{\scriptsize 1084--1084} & A & 1110101 & $\checkmark$ & 274.1 & 47.7 & 174.8 & 17.1 & 4.4 & \hfill 10.3 & \raggedright \citet{Fairchild2017a} \tabularnewline
\raggedright Laurentia & \raggedright Lake Shore Traps & \raggedright 1086\newline{\scriptsize 1084--1087} & A & 1111101 & $\checkmark$ & 272.3 & 47.4 & 185.7 & 23.1 & 3.3 & \hfill 21.8 & \raggedright \citet{Diehl1994a,Kulakov2013a} \tabularnewline
\raggedright Laurentia & \raggedright Schroeder Lutsen Basalts & \raggedright 1090\newline{\scriptsize 1084--1092} & A & 1110101 & $\checkmark$ & 269.1 & 47.5 & 187.8 & 27.1 & 3.0 & \hfill 25.9 & \raggedright \citet{Fairchild2017a} \tabularnewline
\raggedright Laurentia & \raggedright Portage Lake Volcanics & \raggedright 1092\newline{\scriptsize 1092--1093} & A & 1111101 & $\checkmark$ & 271.6 & 47.3 & 183.0 & 27.1 & 2.4 & \hfill 22.7 & \raggedright \citet{Swanson-Hysell2019a} \tabularnewline
\raggedright Laurentia & \raggedright Uppermost Mamainse Point volcanics -N & \raggedright 1094\newline{\scriptsize 1090--1100} & A & 1111111 & $\checkmark$ & 275.2 & 47.1 & 183.2 & 31.2 & 2.5 & \hfill 25.6 & \raggedright \citet{Swanson-Hysell2014a} \tabularnewline
\raggedright Laurentia & \raggedright North Shore Volcanic Group -N (combined) & \raggedright 1095\newline{\scriptsize 1092--1098} & A & 1110101 & $\checkmark$ & 268.7 & 46.3 & 180.1 & 35.1 & 1.9 & \hfill 26.1 & \raggedright \citet{Swanson-Hysell2019a} \tabularnewline
\raggedright Laurentia & \raggedright Central Arizona diabases & \raggedright 1098\newline{\scriptsize 1082--1110} & A & 1111101 &  & 249.5 & 33.8 & 177.4 & 30.6 & 8.9 & \hfill 21.5 & \raggedright \citet{Harlan1993a,Donadini2011b} \tabularnewline
\raggedright Laurentia & \raggedright Mamainse Point volcanics -C (lower N, upper R) & \raggedright 1100\newline{\scriptsize 1100--1101} & A & 1111111 & $\checkmark$ & 275.3 & 47.1 & 189.7 & 36.1 & 4.9 & \hfill 32.9 & \raggedright \citet{Swanson-Hysell2014a} \tabularnewline
\raggedright Laurentia & \raggedright Lower Mamainse Point volcanics -R2 & \raggedright 1105\newline{\scriptsize 1100--1109} & A & 1111111 & $\checkmark$ & 275.3 & 47.1 & 205.2 & 37.5 & 4.5 & \hfill 43.9 & \raggedright \citet{Swanson-Hysell2014a} \tabularnewline
\raggedright Laurentia & \raggedright Osler Volcanic Group reverse upper & \raggedright 1105\newline{\scriptsize 1105--1105} & A & 1111101 & $\checkmark$ & 271.8 & 48.6 & 203.4 & 42.3 & 3.7 & \hfill 45.1 & \raggedright \citet{Swanson-Hysell2019a} \tabularnewline
\raggedright Laurentia & \raggedright Coldwell Complex & \raggedright 1107\newline{\scriptsize 1105--1109} & B & 1110111 & $\checkmark$ & 273.5 & 48.8 & 206.5 & 47.2 & 4.8 & \hfill 49.2 & \raggedright \citet{Kulakov2014c} \tabularnewline
\raggedright Laurentia & \raggedright Osler Volcanic Group reverse middle & \raggedright 1107\newline{\scriptsize 1105--1110} & A & 1110101 & $\checkmark$ & 272.3 & 48.8 & 211.3 & 42.7 & 8.2 & \hfill 50.5 & \raggedright \citet{Swanson-Hysell2014b} \tabularnewline
\raggedright Laurentia & \raggedright Nipigon sills and lavas & \raggedright 1108\newline{\scriptsize 1106--1110} & A & 1011101 & $\checkmark$ & 270.9 & 49.1 & 217.8 & 47.2 & 4.0 & \hfill 56.4 & \raggedright \citet{Evans2021a} \tabularnewline
\raggedright Laurentia & \raggedright Osler Volcanic Group reverse lower & \raggedright 1108\newline{\scriptsize 1105--1110} & A & 1110101 & $\checkmark$ & 272.3 & 48.8 & 218.6 & 40.9 & 4.8 & \hfill 54.6 & \raggedright \citet{Swanson-Hysell2014b} \tabularnewline
\raggedright Laurentia & \raggedright Lowermost Mamainse Point volcanics -R1 & \raggedright 1109\newline{\scriptsize 1106--1112} & A & 1111101 & $\checkmark$ & 275.3 & 47.1 & 227.0 & 49.5 & 5.3 & \hfill 62.9 & \raggedright \citet{Swanson-Hysell2014a} \tabularnewline
\raggedright Laurentia & \raggedright Abitibi Dykes & \raggedright 1141\newline{\scriptsize 1139--1143} & A & 1011111 & $\checkmark$ & 279.0 & 48.0 & 213.4 & 50.4 & 12.6 & \hfill 54.6 & \raggedright \citet{Ernst1993a} \tabularnewline
\raggedright Laurentia & \raggedright NW Ontario Lamprophyre Dykes & \raggedright 1144\newline{\scriptsize 1137--1151} & A & 1111111 & $\checkmark$ & 273.6 & 48.4 & 223.3 & 57.9 & 9.2 & \hfill 61.2 & \raggedright \citet{Piispa2018a} \tabularnewline
\raggedright Laurentia-Greenland & \raggedright NE-SW Trending Dyke Swarm & \raggedright 1160\newline{\scriptsize 1155--1165} & B & 1110101 &  & 314.6 & 61.2 & 231.3 & 34.1 & 5.8 & \hfill 54.2 & \raggedright \citet{Piper1992a} \tabularnewline
\raggedright Laurentia-Greenland & \raggedright Giant Gabbro Dikes & \raggedright 1163\newline{\scriptsize 1161--1165} & B & 1010001 &  & 313.6 & 60.9 & 227.6 & 42.3 & 9.2 & \hfill 56.4 & \raggedright \citet{Piper1977a} \tabularnewline
\raggedright Laurentia-Greenland & \raggedright South Qoroq Intrusion & \raggedright 1163\newline{\scriptsize 1161--1165} & A & 1111001 &  & 314.6 & 61.2 & 217.0 & 42.2 & 14.9 & \hfill 49.6 & \raggedright \citet{Piper1992a} \tabularnewline
\raggedright Laurentia-Greenland & \raggedright Hviddal & \raggedright 1184\newline{\scriptsize 1179--1189} & B & 1010001 &  & 313.7 & 60.9 & 215.6 & 34.1 & 9.6 & \hfill 44.1 & \raggedright \citet{Piper1977a} \tabularnewline
\raggedright Laurentia-Greenland & \raggedright Narsaqq & \raggedright 1184\newline{\scriptsize 1179--1189} & B & 1010001 &  & 313.9 & 60.9 & 225.7 & 31.8 & 10.6 & \hfill 49.2 & \raggedright \citet{Piper1977a} \tabularnewline
\raggedright Laurentia-Scotland & \raggedright Stoer Group & \raggedright 1199\newline{\scriptsize 1129--1269} & B & 0111001 &  & 354.5 & 58.0 & 238.4 & 37.2 & 7.7 & \hfill 43.9 & \raggedright \citet{Evans2021a} \tabularnewline
\raggedright Laurentia & \raggedright Sudbury Dike Swarm & \raggedright 1235\newline{\scriptsize 1232--1242} & A & 1111101 & $\checkmark$ & 278.6 & 46.3 & 192.3 & $-$2.5 & 2.5 & \hfill 8.0 & \raggedright \citet{Palmer1977a} \tabularnewline
\raggedright Laurentia & \raggedright Mackenzie dykes recompiled & \raggedright 1267\newline{\scriptsize 1265--1269} & A & 1111101 & $\checkmark$ & 250.0 & 65.0 & 189.5 & 2.6 & 2.8 & \hfill 9.8 & \raggedright \citet{Fahrig1969a,Robertson1969a,Irving1972c,Park1974a} \tabularnewline
\raggedright Laurentia-Greenland & \raggedright West Gardar Dolerite Dykes & \raggedright 1271\newline{\scriptsize 1263--1278} & B & 0110001 &  & 311.7 & 61.2 & 201.9 & 8.6 & 5.9 & \hfill 17.1 & \raggedright \citet{Piper1977b} \tabularnewline
\raggedright Laurentia-Greenland & \raggedright West Gardar Lamprophyre Dykes & \raggedright 1271\newline{\scriptsize 1263--1278} & B & 0010001 &  & 311.7 & 61.2 & 206.5 & 3.2 & 5.5 & \hfill 15.9 & \raggedright \citet{Piper1977b} \tabularnewline
\raggedright Laurentia-Greenland & \raggedright Kungnat Ring Dyke & \raggedright 1275\newline{\scriptsize 1273--1277} & B & 1010101 &  & 311.7 & 61.2 & 198.6 & 3.4 & 3.3 & \hfill 11.0 & \raggedright \citet{Piper1977b} \tabularnewline
\raggedright Laurentia-Greenland & \raggedright North Qoroq Intrusion & \raggedright 1275\newline{\scriptsize 1274--1276} & B & 1111011 &  & 314.6 & 61.2 & 202.5 & 13.5 & 10.6 & \hfill 21.2 & \raggedright \citet{Piper1992a} \tabularnewline
\raggedright Laurentia & \raggedright Nain Anorthosite & \raggedright 1305\newline{\scriptsize 1283--1328} & B & 0010111 &  & 298.2 & 56.5 & 206.8 & 11.8 & 2.5 & \hfill 28.2 & \raggedright \citet{Murthy1978a} \tabularnewline
\raggedright Laurentia-Greenland & \raggedright Midsommersoe Dolerites & \raggedright 1382\newline{\scriptsize 1380--1384} & B & 1010101 & $\checkmark$ & 333.4 & 81.6 & 242.0 & 10.0 & 3.5 & \hfill 41.7 & \raggedright \citet{Marcussen1983a} \tabularnewline
\raggedright Laurentia-Greenland & \raggedright Victoria Fjord dolerite dykes & \raggedright 1382\newline{\scriptsize 1380--1384} & B & 1011001 & $\checkmark$ & 315.3 & 81.5 & 231.3 & 12.7 & 5.1 & \hfill 38.4 & \raggedright \citet{Abrahamsen1987a} \tabularnewline
\raggedright Laurentia-Greenland & \raggedright Zig-Zag Dal Basalts & \raggedright 1382\newline{\scriptsize 1380--1384} & B & 1010101 & $\checkmark$ & 334.8 & 81.2 & 242.8 & 12.2 & 3.9 & \hfill 44.0 & \raggedright \citet{Marcussen1983a} \tabularnewline
\raggedright Laurentia & \raggedright Pilcher, Garnet Range, Libby & \raggedright 1385\newline{\scriptsize 1378--1407} & B & 1001001 &  & 246.3 & 46.7 & 217.6 & $-$20.6 & 7.9 & \hfill 8.8 & \raggedright \citet{Elston2002a} \tabularnewline
\raggedright Laurentia & \raggedright McNamara & \raggedright 1392\newline{\scriptsize 1378--1407} & B & 1001001 &  & 246.3 & 47.1 & 210.2 & $-$14.3 & 7.9 & \hfill 10.1 & \raggedright \citet{Elston2002a} \tabularnewline
\raggedright Laurentia & \raggedright Purcell Lava & \raggedright 1427\newline{\scriptsize 1412--1442} & B & 1101101 &  & 245.1 & 49.2 & 215.5 & $-$13.6 & 8.2 & \hfill 13.6 & \raggedright \citet{Elston2002a} \tabularnewline
\raggedright Laurentia & \raggedright Mean Rocky Mountain intrusions & \raggedright 1430\newline{\scriptsize 1415--1445} & B & 0110001 & $\checkmark$ & 253.8 & 40.3 & 217.9 & $-$10.4 & 5.9 & \hfill 17.5 & \raggedright \citet{Harlan1994a,Harlan1998a} \tabularnewline
\raggedright Laurentia & \raggedright Mistastin Batholith & \raggedright 1441\newline{\scriptsize 1419--1463} & B & 0011001 & $\checkmark$ & 296.4 & 55.6 & 205.1 & $-$1.4 & 7.3 & \hfill 17.2 & \raggedright \citet{Fahrig1976a,Herve2015a} \tabularnewline
\raggedright Laurentia & \raggedright Snowslip & \raggedright 1449\newline{\scriptsize 1436--1463} & B & 1001001 &  & 245.9 & 47.9 & 212.1 & $-$25.2 & 3.3 & \hfill 2.2 & \raggedright \citet{Elston2002a} \tabularnewline
\raggedright Laurentia & \raggedright Spokane & \raggedright 1458\newline{\scriptsize 1445--1471} & B & 1001001 &  & 246.6 & 48.2 & 216.1 & $-$16.4 & 7.3 & \hfill 11.6 & \raggedright \citet{Elston2002a} \tabularnewline
\raggedright Laurentia & \raggedright St. Francois Mountains Acidic Rocks & \raggedright 1466\newline{\scriptsize 1448--1484} & A & 0111101 & $\checkmark$ & 269.5 & 37.5 & 219.1 & $-$12.1 & 6.0 & \hfill 16.8 & \raggedright \citet{Meert2002b} \tabularnewline
\raggedright Laurentia & \raggedright Michikamau Intrusion Combined & \raggedright 1469\newline{\scriptsize 1468--1470} & A & 1011101 & $\checkmark$ & 296.0 & 54.5 & 217.8 & $-$1.4 & 4.3 & \hfill 24.9 & \raggedright \citet{Emslie1976a} \tabularnewline
\raggedright Laurentia & \raggedright Western Channel diabase & \raggedright 1592\newline{\scriptsize 1589--1595} & A & 1011101 & $\checkmark$ & 242.2 & 66.3 & 245.5 & 9.2 & 4.3 & \hfill 47.9 & \raggedright \citet{Irving1972b} \tabularnewline
\raggedright Laurentia-Greenland & \raggedright Melville Bugt diabase dykes & \raggedright 1630\newline{\scriptsize 1619--1638} & B & 1110101 &  & 301.6 & 74.4 & 274.5 & 4.0 & 9.3 & \hfill 44.5 & \raggedright \citet{Halls2011a} \tabularnewline
\raggedright Laurentia & \raggedright Cleaver Dykes & \raggedright 1740\newline{\scriptsize 1736--1745} & A & 1111101 & $\checkmark$ & 242.0 & 65.7 & 274.9 & 20.4 & 7.0 & \hfill 63.0 & \raggedright \citet{Irving2004a} \tabularnewline
\raggedright Laurentia & \raggedright Wharton Group & \raggedright 1756\newline{\scriptsize 1753--1759} & A & 1011111 & $\checkmark$ & 261.5 & 63.6 & 277.0 & 9.2 & 11.4 & \hfill 51.6 & \raggedright \citet{Raub2026a} \tabularnewline
\raggedright Laurentia & \raggedright NE trending ECMB Diabase Dykes & \raggedright 1779\newline{\scriptsize 1777--1781} & A & 1111101 & $\checkmark$ & 265.8 & 45.5 & 265.8 & 20.5 & 4.5 & \hfill 63.7 & \raggedright \citet{Swanson-Hysell2021b} \tabularnewline
\end{longtable}
\endgroup
